% arara: clean: {
% arara: --> extensions:
% arara: --> ['aux', 'log', 'out', 'pdf']
% arara: --> }
% arara: lualatex: {
% arara: --> shell: no,
% arara: --> draft: yes,
% arara: --> interaction: batchmode
% arara: --> }
% arara: lualatex: {
% arara: --> shell: yes,
% arara: --> draft: no,
% arara: --> interaction: batchmode
% arara: --> }
% arara: clean: {
% arara: --> extensions:
% arara: --> ['aux', 'log', 'out']
% arara: --> }
\documentclass[
	spanish,
	addpoints,
	noanswers,
	a4paper,
	8pt,
	leqno
]{exam}
\usepackage[T1]{fontenc}
\usepackage[spanish]{babel}
\usepackage{libertine}
\usepackage{geometry}
\geometry{
    left = 2cm,
    right = 2cm,
    top = 2cm,
    bottom = 3cm
}
\usepackage[intlimits]{mathtools}
\usepackage{amssymb}
\usepackage{diffcoeff}
\usepackage{amsthm}
\usepackage[svgnames]{xcolor}
\usepackage{siunitx}
\usepackage{tikz}
\usepackage[shortlabels]{enumitem}
\usepackage{multicol}
\usepackage{hyperref}
\theoremstyle{definition}
\newtheorem{definition}{Definición}
\pagestyle{headandfoot}

\usepackage{minted}
\renewcommand{\solutiontitle}{\noindent\textbf{Solución}\par\noindent}
\renewcommand\listingscaption{Listado}
\newcommand{\unmarkedfntext}[1]{%
    \begingroup
    \renewcommand\thefootnote{}\footnote{#1}%
    \addtocounter{footnote}{-1}%
    \endgroup
}

\difdef{c}{L}{op-symbol=\mathop{}\!\mathbin\bigtriangleup}

\pointpoints{Punto}{Puntos}
\bracketedpoints

\difdef {f,fp,s,sp} {|}{
    outer-Ldelim =\left.,
    outer-Rdelim=\right|,
    sub-nudge=0mu
}
\footer{}{Página \thepage\ de \numpages}{}


\begin{document}
\begin{center}
	\sffamily\bfseries\Large
	Métodos Numéricos para Ecuaciones Diferenciales Parciales~CM 032 A\\
	Prueba de entrada
\end{center}

\begin{center}
	\fbox{\fbox{\parbox{5.5in}
			{\centering
				Responda las preguntas en los espacios provistos en
				las hojas de preguntas.
				Los argumentos y la claridad de las respuestas se
				considerarán en la puntuación final.
			}}}
\end{center}
\vspace{0.1in}
\makebox[\textwidth]{Nombres y apellidos del alumno:\enspace\hrulefill}
\vspace{0.2in}
\makebox[\textwidth]{Nombres y apellidos del profesor:\enspace Fidel Jara Huanca.\hfill}

\begin{questions}
	% \header{CM 032}{Prueba de Entrada}{23 de marzo de 2026}
	\question

	Una barra metálica de \qty{100}{\centi\metre} de longitud tiene sus
extremos $x=0$, $x=100$ mantenidos a \qty{0}{\degreeCelsius}.
Inicialmente, la mitad de la barra está a \qty{60}{\degreeCelsius};
mientras que la otra mitad a \qty{40}{\degreeCelsius}.
Asumiendo un coeficiente de difusividad de \num{0.16} unidades c.g.s
y un entorno aislado.

\begin{parts}
	\part\label{1a}

	Describa la ecuación diferencial correspondiente, así como las
	condiciones iniciales y de frontera.

	\part

	No resuelva, solo comente los pasos a tener en cuenta para resolver
	el problema en el ítem~(\ref{1a}).
\end{parts}

\begin{figure}[ht!]
	\centering
	\begin{tikzpicture}[xscale=0.75,yscale=0.75]
		\draw[help lines,color=cyan,thick,densely dotted] (0,0) grid (24,22);
		\draw (1.4,21.4) node {\textsf{\bfseries Solución:}};
	\end{tikzpicture}
\end{figure}


	\question

	Dado el problema
\begin{equation*}
	\begin{cases}\displaystyle
		\diffp{u}{t}=
		\diffp[2]{u}{x}+
		a\diffp{u}{x} &
		\text{en }\mathbb{R}\times\left(0,\infty\right). \\
		u=f           &
		\text{en }\mathbb{R}\times\left\{t=0\right\}.
	\end{cases}
\end{equation*}
Utilizando la transformada de Fourier, determine la solución en el contexto de su transformada.

\begin{figure}[ht!]
	\centering
	\begin{tikzpicture}[xscale=0.75,yscale=0.75]
		\draw[help lines,color=cyan,thick,densely dotted] (0,0) grid (24,29);
		\draw (1.4,28.4) node {\textsf{\bfseries Solución:}};
	\end{tikzpicture}
\end{figure}


	\question

	\begin{parts}
	\part

	Describa cuatro términos de la serie de Taylor de una función real
	de variable real infinitamente diferenciable en un dominio
	adecuado.

	\part

	Describa cuatro términos de la serie de una función real de dos
	variables reales infinitamente diferenciable en un dominio
	adecuado.

	\part

	Describa el gradiente y el hessiano de una función real de tres
	variables reales en un dominio adecuado.
\end{parts}

\begin{figure}[ht!]
	\centering
	\begin{tikzpicture}[xscale=0.75,yscale=0.75]
		\draw[help lines,color=cyan,thick,densely dotted] (0,0) grid (24,26);
		\draw (1.4,25.4) node {\textsf{\bfseries Solución:}};
	\end{tikzpicture}
\end{figure}


\end{questions}
\vfill
\begin{flushright}\bfseries
	Facultad de Ciencias\\[2mm]
	Universidad Nacional de Ingeniería\\[2mm]
	23 de marzo de 2026%\unmarkedfntext{Cada pregunta vale 5 puntos.}
\end{flushright}
\end{document}

